\chapter*{Zusammenfassung}
\thispagestyle{empty}
Die autonome Landung von Drohnen in hindernisreichen Umgebungen erfordert die gleichzeitige Bewältigung von Zielnavigation und Hindernisvermeidung. Bestehende Ansätze behandeln diese Teilaufgaben überwiegend getrennt, und insbesondere in natürlichen Umgebungen mit dichter Vegetation fehlen integrierte Lösungen. Die vorliegende Arbeit untersucht, inwiefern ein Reinforcement-Learning-basierter Ansatz eine autonome Drohne befähigen kann, mithilfe einer minimalen Sensorkonfiguration aus Tiefenkamera und propriozeptivem Zustandsvektor in solchen Umgebungen sicher zu landen. Dazu wird ein Agent mit einer CNN-MLP-Architektur mittels Proximal Policy Optimization in einer simulierten Korridorumgebung mit 3D-Hindernissen im Genesis-Physiksimulator trainiert. Ein kaskadierter PID-Regler überführt die gelernten Positionsbefehle in Motorsteuerungssignale. Die Evaluation erfolgt zweiphasig: In Phase~1 erreicht der Agent in der Trainingsumgebung mit ungesehenen Hindernissen eine Erfolgsrate von 98,4\,\%. In Phase~2 wird die gelernte Strategie im Zero-Shot-Transfer auf eine photogrammetrische Weinbergumgebung übertragen. Hier zeigt sich eine deutliche Leistungsreduktion (10,4\,\% Erfolgsrate bei 0,5\,m Landeradius), wobei der Agent das Ziel zuverlässig anfliegt, jedoch an der terminalen Präzisionslandung scheitert. Das Training erweist sich als instabil: Nur vier von sechs Seeds konvergieren, wobei drei davon nach Erreichen des Optimums einen abrupten Strategiekollaps zeigen, was ein gezieltes Early Stopping erfordert. Die Ergebnisse belegen, dass RL mit minimaler Sensorik die kombinierte Aufgabe in der Trainingsumgebung löst. Der Transfer auf realistische Vegetation ist eingeschränkt möglich, wobei die Leistungsgrenzen in der Trainingsdiversität und fehlender zeitlicher Integration der Beobachtungen liegen.

\emph{Schlüsselwörter:} \\
Reinforcement Learning, autonome Drohnenlandung, Hindernisvermeidung, Tiefensensorik, Proximal Policy Optimization, Genesis-Simulator